\section{Análisis de riesgos}

\subsection{R01. Interrupción de la carga de documentos}

\textbf{Origen.} En el caso Fastly 2021, estudiado en el FORM-AC-01, una configuración válida activó un defecto latente. En el sistema de pre-sustentaciones, un archivo que cumple las reglas de carga podría activar un defecto del servicio e impedir el envío de la solicitud.

\textbf{Probabilidad estimada: media.} El riesgo puede aparecer al procesar archivos con tamaños cercanos al límite permitido.

\textbf{Impacto: alto.} Si el fallo ocurre cerca del cierre, el estudiante podría no completar la entrega dentro del plazo.

\textbf{Mitigación.} Probar archivos por debajo, en el límite y por encima del tamaño permitido. Si la carga falla, mantener la solicitud en borrador, mostrar el error y permitir un reintento. Verificar que el sistema solo confirma el envío cuando el documento queda almacenado.

\textbf{Fuente del caso.} \href{https://www.fastly.com/blog/summary-of-june-8-outage}{Resumen del incidente del 8 de junio de 2021, publicado por Fastly}.

\subsection{R02. Acta sin la nota registrada}

\textbf{Origen.} Aporte de Mario, basado en el incidente de GitHub del 21 de octubre de 2018. En el sistema de pre-sustentaciones, una réplica atrasada podría devolver datos anteriores al último registro del tribunal.

\textbf{Probabilidad estimada: media.} El riesgo aparece si se genera el acta inmediatamente después de guardar la nota y la réplica todavía no ha recibido ese cambio.

\textbf{Impacto: alto.} El tribunal podría firmar un acta sin la nota recién registrada y tendría que corregirla y firmarla de nuevo.

\textbf{Mitigación.} Consultar las notas y generar las actas siempre desde la base de datos principal. Como prueba, registrar una nota, generar el acta de inmediato y comprobar que contiene exactamente la nota guardada.

\textbf{Fuente del caso.} \href{https://github.blog/news-insights/company-news/oct21-post-incident-analysis/}{Análisis del incidente del 21 de octubre, publicado por GitHub}.
